\documentclass[10pt]{article}
\usepackage{neuripslike}
\usepackage{booktabs}
\usepackage{amsmath}
\usepackage{graphicx}
\usepackage{url}
\usepackage[hidelinks]{hyperref}
\usepackage{enumitem}
\setlist{nosep,leftmargin=1.2em}

\newcommand{\placeholder}[1]{{\color{red}\textbf{[#1]}}}

\title{Limitations of Automated Binary Vulnerability Analysis\\in Optimized and Stripped Binaries}
\author{Los Escalante \\ \texttt{losesc@vt.edu} \\ Milestone Report}
\date{}

\begin{document}
\maketitle

\section{Revised scope and thesis}
The proposal asked whether automated static analysis and LLM-based vulnerability
detectors stay reliable once binaries are compiled with aggressive optimization
and stripped of symbols. Three changes were made after instructor feedback and a
closer read of the literature.

First, the baseline changed. The proposal named \texttt{cppcheck}, but
\texttt{cppcheck} analyzes C source and never inspects a compiled binary, so it
cannot measure any effect of optimization or stripping. It is kept only as a
source-level reference point; the binary-level static baseline is now
\texttt{cwe\_checker}, which operates on the lifted binary and covers the three
target weakness classes. Second, the targets are reproduction-grade:
each program has a vulnerable and a fixed (patched) version, so false positives
and hallucination can be measured, not only missed bugs. Third, optimization and
symbol stripping are now varied \emph{independently} in a $2\times2$ design, as
the feedback requested, rather than bundled into a single \texttt{-O3 -s} build.

\textbf{Thesis (current).} We reproduce a published source-level
vulnerability-detection result and test how far it degrades when the same
programs are presented as decompiled binaries, isolating the effect of compiler
optimization (\texttt{-O0} vs.\ \texttt{-O3}) from the effect of symbol stripping
(kept vs.\ stripped), in order to mark the point at which automated triage is no
longer reliable and manual Ghidra/GDB analysis becomes mandatory.

\section{Reproduction target and baseline}
We reproduce a specific result from Khare et al.~\cite{khare2024}, who evaluated
16 LLMs on five vulnerability datasets. We adopt their \emph{Juliet C/C++}
setting and their exact prompts (system prompt, Basic, and CWE-specific) at
temperature~0. The claim we reproduce is their reported Juliet C/C++ detection
performance, summarized in Table~\ref{tab:baseline}. Their best model on this
dataset (Qwen-2.5-14B, CWE-specific prompt) reaches accuracy~0.65 and F1~0.72,
while the CodeQL static baseline reaches accuracy~0.72 but recall only~0.44.
A notable pattern is that most models achieve high recall ($\sim$0.9--1.0) with
precision near~0.5: they call almost everything vulnerable. Our fixed-version
programs are designed to expose exactly this bias and to track whether it worsens
on binaries. Our eight targets follow the Juliet style but are hand-written for this milestone; a Juliet C/C++ subset is integrated (\texttt{fetch\_juliet.sh}) and will replace them for the final report.

\begin{table}[h]
\centering
\small
\begin{tabular}{llcccc}
\toprule
Source (Juliet C/C++) & Method & Acc & Prec & Rec & F1 \\
\midrule
Khare et al.~\cite{khare2024} & CodeQL (static) & 0.72 & 0.99 & 0.44 & 0.61 \\
Khare et al.~\cite{khare2024} & Qwen-2.5-14B (CWE) & 0.65 & 0.60 & 0.89 & 0.72 \\
Khare et al.~\cite{khare2024} & GPT-4 (CWE) & 0.58 & 0.54 & 0.95 & 0.69 \\
\midrule
This work (reproduction) & this work (8 hand-written, source) & 0.875 & 0.875 & 0.875 & 0.875 \\
\bottomrule
\end{tabular}
\caption{Published Juliet C/C++ numbers we reproduce, and our source-level
reproduction row (our 8 hand-written targets, source level).}
\label{tab:baseline}
\end{table}

\section{What is implemented}
The full pipeline runs end to end; only the model inference step requires the
local GPU and is run on the author's machine (Ollama, \texttt{qwen2.5-coder:14b},
temperature~0). The stages are:
\begin{itemize}
\item \textbf{Targets.} Eight C programs, each with a vulnerable and a fixed
variant (16 source cases), covering buffer overflow (CWE-121/787), format string
(CWE-134), and use-after-free / double-free (CWE-416/415). Juliet-style
identifiers are renamed to neutral names so the model cannot read the answer from
a symbol like \texttt{CWE121\_..\_bad}.
\item \textbf{Build matrix.} \texttt{compile.sh} produces the $2\times2$ builds per
variant. Stripping is applied to a \emph{copy} of the compiled binary
(\texttt{strip --strip-all}), not via recompilation, so the machine code is
byte-identical to its symbol-bearing twin. \texttt{\_FORTIFY\_SOURCE} and the
stack protector are disabled so \texttt{-O3} does not replace \texttt{strcpy}
with \texttt{\_\_strcpy\_chk} or insert canaries, which would change the bug
rather than merely optimize around it.
\item \textbf{Build verification.} \texttt{verify\_builds.sh} confirms (1) each
stripped binary has a \texttt{.text} section byte-identical to its kept twin, so
symbols are the only difference, and (2) every vulnerable binary still contains
its bug after compilation --- memory bugs confirmed by AddressSanitizer, format
strings confirmed by a crash under a format-specifier payload --- at both
\texttt{-O0} and \texttt{-O3}. All 16 vulnerable builds are confirmed.
\item \textbf{Decompilation.} \texttt{decompile.sh} exports Ghidra headless
pseudo-C per binary (\texttt{objdump -d} \texttt{.text} fallback when Ghidra is
absent), plus a comment-stripped source corpus for the reproduction condition.
\item \textbf{Inference + scoring.} \texttt{run\_pipeline.py} applies the Khare
prompts and logs a structured verdict per item; \texttt{score.py} computes
accuracy, precision, recall, F1, CWE-match, and hallucination per condition, and
the factor-isolation deltas.
\end{itemize}

\section{Evaluation design and the development/evaluation split}
The detection task is binary (vulnerable vs.\ not), scored against the known
label of each variant. Measures: accuracy, precision, recall, and F1 for
detection; \emph{CWE-match} (of correctly flagged vulnerable cases, the fraction
assigned the right CWE); and \emph{hallucination} (fraction of fixed programs
flagged vulnerable). The five conditions are \texttt{source}, and the four binary
builds \texttt{O0/kept}, \texttt{O0/stripped}, \texttt{O3/kept},
\texttt{O3/stripped}. Because the project reproduces and diagnoses rather than
trains, the ``split'' is a prompt-development vs.\ evaluation split: prompts and
parsing were fixed on two held-out programs (one overflow, one use-after-free)
and never changed based on the remaining evaluation targets. Factor isolation
compares conditions that differ in exactly one variable: \texttt{O0/kept} vs.\
\texttt{O3/kept} isolates optimization; \texttt{O0/kept} vs.\ \texttt{O0/stripped}
isolates stripping.

\section{Preliminary results and interpretation}
The pipeline, baseline reproduction, and verification are complete.
Table~\ref{tab:results} reports detection and triage metrics per condition from
the model run (\texttt{qwen2.5-coder:14b}, temperature~0).

\begin{table}[h]
\centering
\small
\begin{tabular}{lcccccc}
\toprule
Condition & Acc & Prec & Rec & F1 & CWE-match & Halluc. \\
\midrule
source & 0.875 & 0.875 & 0.875 & 0.875 & 1.000 & 0.125 \\
O0 / kept & 0.625 & 0.667 & 0.500 & 0.571 & 1.000 & 0.250 \\
O0 / stripped & 0.562 & 0.667 & 0.250 & 0.364 & 1.000 & 0.125 \\
O3 / kept & 0.600 & 0.600 & 0.429 & 0.500 & 1.000 & 0.250 \\
O3 / stripped & 0.600 & 0.600 & 0.429 & 0.500 & 1.000 & 0.250 \\
\bottomrule
\end{tabular}
\caption{Detection and triage metrics per condition (16 cases each).}
\label{tab:results}
\end{table}

\textbf{Findings.} Detection degrades sharply from source to binary: F1 falls
from 0.875 on source code to 0.571 on the symbol-bearing \texttt{-O0} binary (a
0.30 drop) before any optimization or stripping is applied --- presenting
compiled code instead of source is the single largest effect. Among the binary
factors, symbol stripping hurts more than optimization: at \texttt{-O0},
stripping costs 0.207 F1 (0.571 to 0.364), while \texttt{-O0} to \texttt{-O3}
with symbols kept costs only 0.071. The failure mode is missed detections, not
false alarms --- recall collapses from 0.875 (source) to 0.25
(\texttt{-O0}/stripped) while precision stays near 0.6--0.67, so the model goes
blind rather than noisy. When it does flag a bug it names the correct CWE every
time (CWE-match 1.0 across all conditions), so the breakdown is in detection,
not classification.

\textbf{Class-dependent degradation.} The aggregate hides a strong per-class
split. Use-after-free and double-free (CWE-416/415) are still detected on
\texttt{-O3}-stripped binaries, because the \texttt{malloc}/\texttt{free} call
pattern survives compilation as recognizable PLT calls. Buffer overflows and
format-string bugs (CWE-121/134) are missed on nearly all binary builds, because
the vulnerable operation (\texttt{strcpy}, \texttt{printf(user)}) compiles away
into generic register and memory instructions with no library-call signature.
Automated binary triage is therefore reliable only for the library-call-based
bug classes in this set.

\textbf{Instruction breakdown at \texttt{-O3}.} For 2 of 16 \texttt{-O3} cases
the model abandoned the required verdict format entirely and instead narrated the
disassembly, once misidentifying the target as glibc's \texttt{fwrite}. These are
logged as unparseable and excluded from the metrics. We treat them as a triage
failure in their own right: at high optimization the input is far enough out of
distribution that the model stops performing the task, which supports routing
\texttt{-O3}-stripped binaries to manual Ghidra/GDB analysis.

\section{What will be studied next (the distribution shift)}
The stress condition is the structural distribution shift from source to
\texttt{-O3}-stripped binary: identical semantics, radically different surface
form. The remaining work is: (1) run \texttt{cwe\_checker} on all four binary
builds as the non-LLM binary baseline; (2) define the triage decision rule
concretely --- an alert is trusted only above a recall threshold with LLM and
\texttt{cwe\_checker} agreement, else it is routed to manual Ghidra/GDB --- and
report how often each condition would be routed to manual analysis; (3) a short
qualitative failure analysis mapping specific compiler transformations (inlining,
\texttt{strcpy} vectorization, dead-store elimination) to specific triage errors.

\section{Completion plan}
\begin{itemize}
\item \textbf{By Fri (milestone):} pipeline + baseline reproduction + verification
done (complete); source and binary runs logged; Table~\ref{tab:results} filled.
\item \textbf{Final report:} add \texttt{cwe\_checker} baseline; triage decision
rule and routing rates; failure analysis; expand targets if time allows.
\item \textbf{Risks:} a bug optimized away at \texttt{-O3} (mitigated by the ASan
survival check, all 16 currently pass); Ghidra headless setup (mitigated by the
\texttt{objdump} fallback, already working).
\end{itemize}

\section*{Class-survey entry}
This project studies binary vulnerability triage: deciding, without source code,
whether a compiled program contains a memory-safety or format-string bug. It
tests whether Khare et al.'s source-level LLM detection results on Juliet C/C++
hold when the model sees only decompiled binaries that are optimized and/or
stripped. The method is a $2\times2$ design (\texttt{-O0}/\texttt{-O3} $\times$
symbols kept/stripped) that scores the same LLM and prompts against the
source-level baseline, with \texttt{cwe\_checker} planned as a non-LLM binary
baseline. The goal is to diagnose where automated triage stops being reliable
and manual Ghidra/GDB analysis becomes necessary.

\begin{center}\small
\begin{tabular}{@{}p{1.6cm}p{2.0cm}p{2.2cm}p{2.6cm}p{1.6cm}p{1.6cm}@{}}
\toprule
Domain & Task & Baseline & Distribution shift studied & Primary metric & Finding so far \\
\midrule
Binary vuln.\ triage & Bug detection + CWE ID (Juliet C/C++) & Khare et al.\ source-level LLM (Qwen-2.5-14B) &
Source $\to$ decompiled binary: \texttt{-O0}/\texttt{-O3} $\times$ kept/stripped & F1 (plus CWE match) & Source F1 0.88 $\to$ 0.36 stripped; stripping $>$ optimization \\
\bottomrule
\end{tabular}
\end{center}

\section*{Contributions}
Solo project. Los Escalante: target curation, build and verification scripts,
decompilation pipeline, prompt pipeline and scorer, reproduction of the Khare
et al.\ baseline, and this report.

\section*{Artifacts}
Code and data: \url{github.com/losesc/binvuln}. Scripts: \texttt{compile.sh},
\texttt{verify\_builds.sh}, \texttt{decompile.sh}, \texttt{run\_pipeline.py},
\texttt{score.py}.

\begin{thebibliography}{9}
\bibitem{khare2024} A.~Khare, S.~Dutta, Z.~Li, A.~Solko-Breslin, R.~Alur, M.~Naik.
``Understanding the Effectiveness of Large Language Models in Detecting Security
Vulnerabilities.'' arXiv:2311.16169, 2024 (ICST 2025).
\bibitem{juliet} T.~Boland, P.~E.~Black. ``Juliet 1.1 C/C++ and Java Test Suite.''
\textit{IEEE Computer}, 2012. NIST SAMATE SARD.
\bibitem{cwechecker} Fraunhofer FKIE. ``cwe\_checker.''
\url{https://github.com/fkie-cad/cwe_checker}.
\end{thebibliography}

\end{document}
